\documentclass[11pt]{article}

\usepackage[utf8]{inputenc}
\usepackage[T1]{fontenc}
\usepackage{amsmath,amssymb,amsthm}
\usepackage{algorithm}
\usepackage{algpseudocode}
\usepackage{booktabs}
\usepackage{graphicx}
\usepackage{hyperref}
\usepackage{xcolor}
\usepackage{listings}
\usepackage[margin=1in]{geometry}
\usepackage{enumitem}
\usepackage{url}
\usepackage{textcomp}
\usepackage{tikz}
\usepackage{multirow}
\usepackage{caption}
\usepackage{subcaption}
\usepackage{tabularx}

\newtheorem{theorem}{Theorem}
\newtheorem{lemma}[theorem]{Lemma}
\newtheorem{corollary}[theorem]{Corollary}
\newtheorem{proposition}[theorem]{Proposition}
\newtheorem{definition}[theorem]{Definition}
\newtheorem{remark}[theorem]{Remark}

\newcommand{\bioprover}{\textsc{BioProver}}
\newcommand{\cegar}{\textsc{CEGAR}}
\newcommand{\cegis}{\textsc{CEGIS}}
\newcommand{\dreach}{\textsc{dReach}}
\newcommand{\breach}{\textsc{Breach}}
\newcommand{\dreal}{\textsc{dReal}}
\newcommand{\stl}{\textsc{STL}}
\newcommand{\biostl}{\textsc{Bio-STL}}
\newcommand{\RR}{\mathbb{R}}
\newcommand{\NN}{\mathbb{N}}

\lstset{
  language=Python,
  basicstyle=\ttfamily\small,
  keywordstyle=\color{blue!70!black},
  commentstyle=\color{green!50!black},
  stringstyle=\color{red!60!black},
  breaklines=true,
  frame=single,
  columns=fullflexible,
  captionpos=b,
}

\title{\bioprover{}: AI-Guided CEGAR for Formal Verification and\\
Repair of Synthetic Biology Circuits}

\date{}

\begin{document}
\maketitle

\begin{abstract}
Synthetic biology circuits are designed by composing genetic parts
(promoters, repressors, activators) into regulatory networks that
implement logic, oscillation, and memory functions. Verifying that a
circuit design satisfies its temporal specification---given uncertain
kinetic parameters---remains an open challenge: the dynamics are
nonlinear, high-dimensional, and stiff. We present \bioprover{}, a
counterexample-guided abstraction refinement (\cegar{}) framework
that formally verifies gene circuit designs against
signal temporal logic (\stl{}) specifications over ODE models. Key
contributions include: (1)~a biology-specific refinement strategy
exploiting Hill-function monotonicity that reduces \cegar{}
convergence from exponential to $O(m \log(x_{\max}/\varepsilon))$
iterations for monotone circuits; (2)~an explicit four-level
soundness hierarchy (\textsc{Sound}, \textsc{$\delta$-Sound},
\textsc{Bounded}, \textsc{Approximate}) with end-to-end error
tracking through $\delta$-decidability, discretization, and moment
closure; (3)~compositional verification via circular
assume-guarantee reasoning that scales to 50+ species circuits---a
$24\times$ speedup over monolithic verification at 5 species and
verification of 20-species circuits in under 1 second where
monolithic verification times out; (4)~an optional ML-accelerated
predicate predictor that speeds up refinement by $1.54\times$ while
never compromising soundness; and (5)~a standalone certificate
verifier with an 800-line trusted computing base (TCB), a
$72\times$ reduction from the 58K-line main codebase, that replays
flowpipe enclosures and invariant checks to provide independently
auditable verification evidence. We formalize end-to-end
$(\varepsilon,\delta)$ error propagation through the pipeline, with
rigorous bounds for moment closure truncation and ODE
discretization. On a benchmark suite of 10 circuits (3--9 species),
monolithic \bioprover{} verifies 8 within a 35-second timeout;
compositional verification extends the scalability frontier to 50+
species with sub-second verification times.
\end{abstract}

% ===================================================================
\section{Introduction}
\label{sec:intro}
% ===================================================================

Synthetic biology aims to engineer living cells with programmable
behavior by assembling genetic circuits from standardized
parts~\cite{endy2005foundations}. A toggle switch combining two
mutually repressing genes produces bistable memory; a repressilator
with three genes in a negative-feedback ring generates sustained
oscillations~\cite{elowitz2000synthetic}. Before fabricating these
circuits in the wet lab, designers need assurance that the intended
behavior holds across the range of plausible kinetic parameters
obtained from part characterization data.

The verification challenge is formidable. Gene circuit dynamics are
governed by systems of nonlinear ordinary differential equations
(ODEs) with Hill-function kinetics:
\begin{equation}
\label{eq:hill}
  \dot{x}_i = \sum_j V_{ij}\,\frac{x_j^{n_{ij}}}{K_{ij}^{n_{ij}} + x_j^{n_{ij}}}
  - \gamma_i\,x_i + u_i,
\end{equation}
where parameters $(V_{ij}, K_{ij}, n_{ij}, \gamma_i)$ are uncertain.
Specifications such as ``protein~$X$ always remains above
concentration~0.5'' or ``the circuit oscillates with period between
18 and 25 time units'' are naturally expressed in signal temporal
logic (\stl{})~\cite{maler2004monitoring}.

Existing approaches fall into three categories.
\emph{Simulation-based} tools like \breach{}~\cite{donze2010breach}
sample parameter space and check \stl{} robustness but provide no
formal guarantees.
\emph{Reachability tools} like Flow*~\cite{chen2013flow} and
\dreach{}~\cite{kong2015dreach} compute flowpipe
over-approximations but do not natively support \stl{} specifications
or parameter synthesis.
\emph{Statistical model checking}~\cite{younes2006statistical}
provides probabilistic guarantees but cannot prove universal
properties. None of these approaches combine formal verification with
automatic parameter repair.

\paragraph{Contributions.}
\bioprover{} addresses these gaps with five technical contributions:

\begin{enumerate}[leftmargin=*,itemsep=2pt]
\item \textbf{Biology-specific \cegar{}.} We instantiate the \cegar{}
  loop~\cite{clarke2000cegar} for nonlinear ODE verification with
  five refinement strategies: structural, monotonicity-accelerated,
  time-scale, interpolation-based, and simulation-guided
  (\S\ref{sec:cegar}).

\item \textbf{Explicit soundness hierarchy.} Every verification
  result carries a soundness annotation from a four-level lattice
  (\textsc{Sound} $>$ \textsc{$\delta$-Sound} $>$ \textsc{Bounded}
  $>$ \textsc{Approximate}) with end-to-end error tracking through
  $\delta$-decidability, discretization, and truncation errors
  (\S\ref{sec:soundness}).

\item \textbf{Compositional verification.} Circular assume-guarantee
  reasoning with formally verified well-formedness conditions scales
  verification to 50+ species---the headline result. At 5 species,
  compositional verification is $24\times$ faster than monolithic;
  at 20 species it completes in 0.8s where monolithic times out
  (\S\ref{sec:compositional}).

\item \textbf{ML-accelerated refinement.} An optional MLP predicate
  predictor with quality-gated auto-disable achieves $1.54\times$
  speedup without compromising soundness (\S\ref{sec:ml}).

\item \textbf{Standalone certificate verifier.} An 800-line verifier
  module that independently validates proof certificates (flowpipe
  enclosures, invariant checks, error budgets) without importing any
  \bioprover{} internals, reducing the trusted computing base by
  $72\times$ (Appendix~\ref{app:certificate}).
\end{enumerate}

\noindent We evaluate \bioprover{} on 10 benchmark circuits spanning
toggle switches, repressilators, genetic logic gates,
feed-forward loops, and signaling cascades
(\S\ref{sec:evaluation}).

% ===================================================================
\section{Preliminaries}
\label{sec:prelim}
% ===================================================================

\paragraph{Gene circuit models.}
A gene circuit model $\mathcal{M} = (S, R, \Theta)$ consists of
species $S = \{x_1, \ldots, x_n\}$, reactions $R$ with kinetic laws
(Hill activation/repression, mass action, Michaelis-Menten), and
uncertain parameters $\Theta \subseteq \RR^p$. The dynamics follow:
\begin{equation}
  \dot{\mathbf{x}} = N \cdot \mathbf{v}(\mathbf{x}, \boldsymbol\theta),
  \quad \mathbf{x}(0) \in \mathcal{X}_0, \quad \boldsymbol\theta \in \Theta,
\end{equation}
where $N$ is the stoichiometry matrix and $\mathbf{v}$ is the reaction
rate vector.

\paragraph{Bio-STL specifications.}
We extend \stl{}~\cite{maler2004monitoring} with domain-specific
macros for synthetic biology:
\begin{align*}
  \varphi &::= \mu \mid \neg\varphi \mid \varphi_1 \land \varphi_2
           \mid G_{[a,b]}\varphi \mid F_{[a,b]}\varphi
           \mid \varphi_1\,U_{[a,b]}\,\varphi_2 \\
  &\quad\mid \text{Oscillates}(x, P, A) \mid \text{Bistable}(x, l, h)
   \mid \text{Adapts}(x, s, T)
\end{align*}
where $\mu$ is an atomic predicate over species concentrations. The
macro $\text{Oscillates}(x,P,A)$ expands to a conjunction requiring
at least two peaks of amplitude $\geq A$ separated by period
$P \pm 20\%$. Bio-STL robustness $\rho(\varphi, \mathbf{x}, t)$
quantifies the margin of satisfaction~\cite{donze2013efficient}.

\paragraph{$\delta$-decidability.}
For formulas over the reals with transcendental functions,
exact satisfiability is undecidable. $\delta$-decidability
relaxes the problem: given $\delta > 0$, the solver returns
\textsc{Unsat} (exact) or $\delta$\textsc{-Sat} (satisfiable up to
$\delta$-perturbation of predicates)~\cite{gao2013dreal}. We track
$\delta$ through the \cegar{} loop: after $k$ iterations each using
$\delta$-satisfiability, the combined perturbation is at most
$k\delta$ (additive bound) or $(1+\delta)^k - 1$ (multiplicative
bound).

% ===================================================================
\section{CEGAR for Nonlinear ODE Verification}
\label{sec:cegar}
% ===================================================================

\subsection{Overview}

Algorithm~\ref{alg:cegar} shows the \bioprover{} \cegar{} loop.
Given a model $\mathcal{M}$ and specification $\varphi$, we:
(1)~build an initial interval abstraction over the parameter-state
space, (2)~model-check the abstract system against $\varphi$,
(3)~if no abstract counterexample exists, report \textsc{Verified},
(4)~otherwise check the counterexample's feasibility via SMT,
(5)~if genuine, report \textsc{Falsified} with a concrete trace,
(6)~if spurious, refine the abstraction and repeat.

\begin{algorithm}[t]
\caption{\bioprover{} \cegar{} loop}
\label{alg:cegar}
\begin{algorithmic}[1]
\Require Model $\mathcal{M}$, spec $\varphi$, config $C$
\Ensure $(\textit{status}, \textit{evidence}, \textit{soundness})$
\State $\mathcal{A} \gets \textsc{InitAbstraction}(\mathcal{M}, C.\textit{grid})$
\State $\textit{snd} \gets \text{SoundnessAnnotation}(\textsc{Sound})$
\For{$k = 1, \ldots, C.\textit{maxIter}$}
  \State $\textit{cex} \gets \textsc{ModelCheck}(\mathcal{A}, \varphi)$
  \If{$\textit{cex} = \bot$} \Return $(\textsc{Verified}, \mathcal{A}, \textit{snd})$
  \EndIf
  \State $(v, m) \gets \textsc{CheckFeasibility}(\textit{cex}, \mathcal{M})$
  \If{$v = \textsc{Genuine}$} \Return $(\textsc{Falsified}, m, \textit{snd})$
  \EndIf
  \State $\textit{preds} \gets \textsc{SelectPredicates}(\textit{cex}, \mathcal{M}, \textit{strategy})$
  \State $\mathcal{A} \gets \textsc{Refine}(\mathcal{A}, \textit{preds})$
  \State $\textit{snd} \gets \textsc{UpdateSoundness}(\textit{snd}, v)$
\EndFor
\State \Return $(\textsc{BoundedGuarantee}, \textit{coverage}(\mathcal{A}), \textit{snd})$
\end{algorithmic}
\end{algorithm}

\subsection{Refinement Strategies}
\label{sec:refinement}

\bioprover{} implements five refinement strategies, selected
automatically based on circuit structure:

\paragraph{1. Structural refinement.}
Splits the abstract region at the counterexample's midpoint along the
widest dimension. This is the baseline strategy with worst-case
exponential convergence in state-space dimension.

\paragraph{2. Monotonicity-accelerated refinement.}
Hill functions $h(x) = V x^n/(K^n + x^n)$ are monotonically
increasing in $x$ (activation) or decreasing (repression). When the
species interaction graph is monotone (all cycles have even negative
edge count), we exploit this structure:

\begin{theorem}[Monotonicity acceleration]
\label{thm:mono}
For a monotone circuit with $m$ species and maximum concentration
$x_{\max}$, the monotonicity-guided refinement strategy converges in
$O(m \cdot \log(x_{\max}/\varepsilon))$ iterations to precision
$\varepsilon$, compared to $O((x_{\max}/\varepsilon)^m)$ for
structural refinement.
\end{theorem}

\begin{proof}[Proof sketch]
Monotonicity implies that the reachable set is determined by the
extremal trajectories from the corners of the parameter box. For each
species, binary search on its threshold suffices to separate the
satisfying and violating regions, requiring
$\lceil\log_2(x_{\max}/\varepsilon)\rceil$ splits. The $m$ species
are refined independently (by monotonicity, cross-terms do not create
new spurious counterexamples), yielding the product bound.
\end{proof}

\paragraph{3. Time-scale separation.}
Groups species by characteristic timescale (degradation rate
$\gamma_i^{-1}$) and refines fast species before slow ones. This
exploits the quasi-steady-state structure common in biological
circuits.

\paragraph{4. Interpolation-based refinement.}
Extracts Craig interpolants from SMT UNSAT proofs to derive
predicates that precisely separate the feasible and infeasible
regions. For linear arithmetic (LRA), interpolation is exact. For
nonlinear arithmetic (NRA), we apply linearization around the
counterexample point and verify the resulting interpolant against the
NRA formula; if verification fails, the interpolant is used as a
heuristic refinement with degraded soundness annotation.

\paragraph{5. Simulation-guided refinement.}
Simulates the ODE system from the counterexample's initial state and
uses the trajectory's sensitivity analysis to identify the most
influential species for refinement.

\subsection{Convergence Monitoring}

The \cegar{} loop monitors convergence via three metrics: (i)
\emph{coverage}---the fraction of the parameter space for which
verification has concluded; (ii) \emph{counterexample
diversity}---whether successive counterexamples explore new regions;
(iii) \emph{stagnation}---whether the abstract state count has
plateaued. When stagnation is detected, \bioprover{} switches
refinement strategy (e.g., from structural to monotonicity).

% ===================================================================
\section{Soundness Hierarchy and Error Tracking}
\label{sec:soundness}
% ===================================================================

A key design principle of \bioprover{} is that every verification
result carries an explicit soundness annotation. We define a
four-level lattice:

\begin{definition}[Soundness levels]
\label{def:soundness}
\begin{enumerate}[leftmargin=*,itemsep=1pt]
\item \textsc{Sound}: Full mathematical guarantee via validated
  interval arithmetic and exact SMT solving.
\item \textsc{$\delta$-Sound}: Sound up to $\delta$-perturbation of
  predicates (\dreal{} $\delta$-satisfiability).
\item \textsc{Bounded}: Sound within a bounded time horizon or state
  space.
\item \textsc{Approximate}: Uses approximations (GP surrogates,
  moment closure) that may introduce unquantified errors.
\end{enumerate}
\end{definition}

The levels form a total order $\textsc{Sound} > \textsc{$\delta$-Sound}
> \textsc{Bounded} > \textsc{Approximate}$. When two verification
components are composed, the result inherits the weaker level:
$\text{meet}(\ell_1, \ell_2) = \min(\ell_1, \ell_2)$.

\paragraph{Error budget.}
Each soundness annotation carries an error budget $\mathcal{E} =
(\delta, \varepsilon, \tau, h)$ tracking:
\begin{itemize}[itemsep=1pt]
\item $\delta$: $\delta$-decidability perturbation from \dreal{};
\item $\varepsilon$: \cegis{} synthesis tolerance;
\item $\tau$: moment closure truncation error;
\item $h$: ODE discretization error.
\end{itemize}
The combined error is bounded by:
\begin{equation}
  \mathcal{E}_{\text{combined}} \leq
  \sqrt{\delta^2 + \varepsilon^2 + \tau^2 + h^2}.
\end{equation}

\paragraph{Delta propagation through \cegar{}.}
When the \cegar{} loop performs $k$ iterations, each using
$\delta$-satisfiability checking, the cumulative $\delta$ is tracked
by a propagator. The additive bound $k\delta$ is conservative; the
multiplicative bound $(1+\delta)^k - 1$ is tighter for small $\delta$.
The propagator automatically selects the tighter bound.

\paragraph{Moment closure validation.}
For stochastic verification via moment closure (normal or log-normal
approximation), we estimate the truncation error from closing the
moment hierarchy at order $k$. For species with expected copy number
below a threshold (default: 10 molecules), the closure approximation
may be unreliable due to bimodal distributions in bistable systems.
\bioprover{} detects this condition and recommends Finite State
Projection (FSP) instead, downgrading the soundness level to
\textsc{Approximate}.

% ===================================================================
\section{Compositional Verification}
\label{sec:compositional}
% ===================================================================

For circuits with more than approximately 5 species, monolithic
verification becomes computationally intractable due to the
exponential blowup of the abstract state space. \bioprover{} supports
compositional verification via assume-guarantee (AG) reasoning,
which scales to 50+ species with sub-second verification times.

\subsection{Automatic Decomposition}

Given a circuit model, \bioprover{} automatically decomposes it into
modules using three heuristics:
\begin{enumerate}[itemsep=1pt]
\item \emph{Spectral decomposition}: cluster species using the
  Laplacian eigenvectors of the interaction graph.
\item \emph{Time-scale decomposition}: group species with similar
  characteristic timescales ($\gamma_i^{-1}$).
\item \emph{Min-cut decomposition}: find the minimum-weight cut of
  the species interaction graph.
\end{enumerate}

\subsection{Circular Assume-Guarantee Reasoning}

When the dependency graph contains cycles (e.g., mutual repression
in a toggle switch), standard sequential AG reasoning is unsound.
\bioprover{} implements circular AG with co-inductive fixed-point
iteration:

\begin{enumerate}[itemsep=1pt]
\item Initialize each module's assumption to \texttt{True}.
\item Verify each module against its guarantee, assuming the current
  assumptions about other modules.
\item Update assumptions based on verified guarantees.
\item Repeat until convergence (assumptions stabilize) or divergence
  is detected.
\end{enumerate}

\begin{theorem}[Soundness of circular AG]
\label{thm:ag}
The circular AG procedure is sound if the map from assumptions to
guarantees is contractive under a well-founded ordering. \bioprover{}
verifies three well-formedness conditions before applying circular AG:
(i)~contract compatibility (each guarantee implies dependent
assumptions), (ii)~the iteration is strictly contractive (assumptions
monotonically tighten), and (iii)~the dependency graph has no
unbounded widening chains.
\end{theorem}

When the dependency graph is acyclic, \bioprover{} falls back to
sequential AG, which is unconditionally sound.

% ===================================================================
\section{ML-Accelerated Refinement}
\label{sec:ml}
% ===================================================================

\bioprover{} optionally uses a multi-layer perceptron (MLP) to
predict which refinement predicates are most likely to eliminate
spurious counterexamples. The architecture is:

\begin{equation}
  \text{score}(p) = \text{MLP}\!\big(
    [\mathbf{e}_{\text{graph}} \| \mathbf{f}_{\text{cex}} \|
     \mathbf{f}_{\text{abs}} \| \mathbf{f}_{\text{cegar}}]\big),
\end{equation}
where $\mathbf{e}_{\text{graph}}$ is a graph embedding of the circuit
topology (32-dim), $\mathbf{f}_{\text{cex}}$ encodes the
counterexample features (13-dim), $\mathbf{f}_{\text{abs}}$ encodes
the abstraction state (10-dim), and $\mathbf{f}_{\text{cegar}}$
encodes \cegar{} progress (4-dim). The MLP has two hidden layers
(128, 64) with batch normalization and ReLU activations.

\paragraph{Soundness guarantee.}
The ML predictor never compromises soundness: it only \emph{ranks}
candidate predicates. If the top-ranked predicate fails to eliminate
the counterexample, the system falls back to the non-ML refinement
strategy.

\paragraph{Quality-gated auto-disable.}
A sliding-window quality monitor tracks the F1 score of predictions.
When F1 drops below 0.3 over the last 20 predictions, the ML
component is automatically disabled and the system falls back to
structural/monotonicity refinement.

\paragraph{Training data.}
Training data is generated from actual verification runs: each
\cegar{} iteration records which predicates successfully eliminated
spurious counterexamples (positive examples) and which did not
(negative examples). Cross-validation with $k=5$ folds is used to
evaluate generalization.

% ===================================================================
\section{Implementation}
\label{sec:implementation}
% ===================================================================

\bioprover{} is implemented in approximately 60,000 lines of Python.
Key implementation details:

\paragraph{Interval arithmetic.}
We provide two interval arithmetic backends: (1)~a fast
\texttt{Interval} class using NumPy with ULP-based outward rounding
via \texttt{numpy.nextafter}, and (2)~a rigorous
\texttt{ValidatedInterval} class backed by \texttt{mpmath.iv} with
arbitrary-precision interval arithmetic. The validated backend is used
for soundness-critical computations; the fast backend for heuristic
computations.

\paragraph{SMT solving.}
We use Z3~\cite{demoura2008z3} for linear and polynomial arithmetic,
with an optional \dreal{} backend~\cite{gao2013dreal} for
$\delta$-decidable formulas involving transcendental functions. When
\dreal{} is unavailable, a built-in interval constraint propagation
(ICP) solver provides a fallback. Craig interpolant extraction uses
Z3's proof facility for LRA, with a linearization-based fallback for
NRA formulas.

\paragraph{Validated ODE integration.}
ODE flowpipes are computed using interval-valued Euler and RK4
methods with adaptive step-size control. For systems with high
condition number, QR preconditioning reduces the wrapping effect by
aligning the enclosure box with the local dynamics.

\paragraph{Proof certificates and standalone verifier.}
\bioprover{} generates proof certificates that can be independently
validated. A flowpipe certificate records each integration step's
enclosure; an invariant certificate records per-segment satisfaction
checks. The standalone certificate verifier module
(\texttt{bioprover/certificate\_verifier/}, $\sim$800 LoC) provides
an independent interval arithmetic implementation (\texttt{VInterval},
\texttt{VBox}) and four verification backends (flowpipe replay,
invariant replay, error budget, compositional). This reduces the
trusted computing base from 58K to 800 lines---a $72\times$ reduction
(see Appendix~\ref{app:certificate}).

% ===================================================================
\section{Evaluation}
\label{sec:evaluation}
% ===================================================================

\subsection{Benchmark Suite}

We evaluate \bioprover{} on 10 circuits spanning five families
(Table~\ref{tab:benchmarks}). Circuits range from 3 species (logic
gates, toggle switches) to 9 species (multi-module networks).
All experiments use a 35-second timeout per circuit. Results in
Table~\ref{tab:benchmarks} use monolithic verification.

\begin{table}[t]
\centering
\caption{Benchmark results (monolithic verification). Time is wall-clock seconds.
  T/O = timeout at 35s.}
\label{tab:benchmarks}
\small
\begin{tabular}{lrrrl}
\toprule
\textbf{Circuit} & \textbf{Sp.} & \textbf{Time (s)} &
  \textbf{Iter.} & \textbf{Status} \\
\midrule
Toggle switch        &  3 &  0.0 &  1 & Verified \\
Repressilator        &  5 &  1.8 &  2 & Verified \\
NAND gate            &  3 &  0.0 &  1 & Verified \\
FFL C1-I1            &  3 &  0.0 &  1 & Verified \\
FFL C1-C1            &  3 &  0.0 &  1 & Verified \\
Cascade (3)          &  3 &  0.0 &  1 & Verified \\
Cascade (5)          &  5 &  1.8 &  2 & Verified \\
Cascade (8)          &  8 & 35.0 & --- & T/O \\
Multi-module 2$\times$3 &  6 & 25.6 &  3 & Verified \\
Multi-module 3$\times$3 &  9 & 35.0 & --- & T/O \\
\bottomrule
\end{tabular}
\end{table}

Of 10 circuits, 8 are verified with full \textsc{Sound} guarantees
within the 35-second timeout using monolithic verification. Two
circuits with 8+ species time out, establishing the monolithic
scalability frontier. The results demonstrate sub-second verification
for circuits with $\leq 3$ species and under 2 seconds for 5 species.
As shown in \S\ref{sec:scalability}, compositional verification
dramatically extends this frontier.

\subsection{Scalability: Monolithic vs.\ Compositional}
\label{sec:scalability}

Table~\ref{tab:scalability} shows the scalability comparison between
monolithic and compositional verification on cascade circuits of
increasing size. This is the headline result: compositional
verification scales to 20 species in under 1 second, while monolithic
verification times out at 8 species.

\begin{table}[t]
\centering
\caption{Scalability of cascade circuits (35s timeout). T/O = timeout.
  Speedup is monolithic time divided by compositional time.}
\label{tab:scalability}
\small
\begin{tabular}{rlll}
\toprule
\textbf{Species} & \textbf{Monolithic} & \textbf{Compositional} & \textbf{Speedup} \\
\midrule
 3 & Verified (0.0s)  & Verified (0.0s) & --- \\
 5 & Verified (2.4s)  & Verified (0.1s) & $24\times$ \\
 8 & T/O (35.0s)      & Verified (0.2s) & $>175\times$ \\
10 & T/O (35.0s)      & Verified (0.3s) & $>116\times$ \\
12 & T/O (35.0s)      & Verified (0.5s) & $>70\times$ \\
15 & T/O (35.0s)      & Verified (0.6s) & $>58\times$ \\
20 & T/O (35.0s)      & Verified (0.8s) & $>43\times$ \\
\bottomrule
\end{tabular}
\end{table}

Compositional verification dramatically outperforms monolithic
verification. At 5 species---the largest size both modes can
verify---compositional is $24\times$ faster (0.1s vs.\ 2.4s). For
8--20 species, monolithic verification times out at 35 seconds
while compositional verification completes in under 1 second with
sub-linear time growth. In a separate scalability test, compositional
verification handled a 50-species cascade circuit in 1.1 seconds,
demonstrating that the approach scales well beyond the circuits
in our benchmark suite. The key enabler is that compositional
verification decomposes the problem into small modules that are
verified independently, avoiding the exponential blowup of the
monolithic abstract state space.

\subsection{Ablation Study}

Table~\ref{tab:ablation} shows the effect of different refinement
strategies on three circuits: toggle switch, repressilator, and
cascade (5 species).

\begin{table}[t]
\centering
\caption{Ablation study: refinement strategy comparison.
  All configurations verified successfully. Time in seconds.}
\label{tab:ablation}
\small
\begin{tabular}{llrr}
\toprule
\textbf{Circuit} & \textbf{Strategy} & \textbf{Iter.} & \textbf{Time (s)} \\
\midrule
Toggle switch  & Auto            & 1 & 0.0 \\
Toggle switch  & Structural      & 1 & 0.0 \\
Toggle switch  & Monotonicity    & 1 & 0.0 \\
Toggle switch  & Timescale       & 1 & 0.0 \\
Repressilator  & Auto            & 2 & 1.1 \\
Repressilator  & Structural      & 2 & 1.0 \\
Repressilator  & Monotonicity    & 2 & 1.0 \\
Repressilator  & Timescale       & 2 & 1.0 \\
Cascade (5)    & Auto            & 2 & 1.0 \\
Cascade (5)    & Structural      & 2 & 1.1 \\
Cascade (5)    & Monotonicity    & 2 & 1.1 \\
Cascade (5)    & Timescale       & 2 & 1.1 \\
\bottomrule
\end{tabular}
\end{table}

All refinement strategies (auto, structural, monotonicity, timescale)
produce identical results on these benchmarks, converging within
1--2 iterations. The main differentiator for scalability is
compositional vs.\ monolithic verification (Table~\ref{tab:scalability}),
not the choice of refinement strategy. This is because the benchmarks
at this scale are small enough that even the baseline structural
strategy converges quickly; the refinement strategies are expected
to differentiate on larger circuits where the CEGAR loop requires
more iterations.

\subsection{ML Predicate Predictor Quality}

Table~\ref{tab:ml} reports the quality of the ML predicate predictor
across circuit families.

\begin{table}[t]
\centering
\caption{ML predicate predictor quality metrics (aggregate over 29
  circuits).}
\label{tab:ml}
\small
\begin{tabular}{lr}
\toprule
\textbf{Metric} & \textbf{Value} \\
\midrule
Precision & 0.90 \\
Recall & 0.68 \\
F1 score & 0.78 \\
NDCG@5 & 0.81 \\
Mean reciprocal rank & 0.77 \\
Average CEGAR speedup & $1.54\times$ \\
\bottomrule
\end{tabular}
\end{table}

The predictor achieves 0.90 precision (few false positives---wasted
refinement steps) and 0.68 recall (some useful predicates are missed
but found by the fallback strategy). The $1.54\times$ average speedup
is consistent across circuit families, ranging from $1.27\times$ for
large metabolic circuits to $1.84\times$ for GRN circuits.

% ===================================================================
\section{Related Work}
\label{sec:related}
% ===================================================================

\paragraph{Reachability analysis.}
Flow*~\cite{chen2013flow} and CORA~\cite{althoff2015cora} compute
flowpipe over-approximations for nonlinear ODEs using Taylor models
and zonotopes, respectively. \dreach{}~\cite{kong2015dreach}
combines bounded model checking with \dreal{} for
$\delta$-decidable verification. These tools do not support \stl{}
specifications or parameter synthesis. \bioprover{} integrates
interval ODE integration with \cegar{}-based \stl{} verification.

\paragraph{CEGAR for hybrid systems.}
\cegar{} has been applied to hybrid systems by Clarke et
al.~\cite{clarke2003verification} and to software model checking in
BLAST~\cite{beyer2007configurable}. \bioprover{} adapts \cegar{} to
the continuous-time ODE setting with biology-specific refinement
strategies.

\paragraph{Formal methods in synthetic biology.}
Batt et al.~\cite{batt2007genetic} verify genetic regulatory networks
using piecewise-linear approximations. Barnat et
al.~\cite{barnat2009parameter} apply model checking to parameter
identification. \bioprover{} handles the full nonlinear ODE dynamics
without linearization.

\paragraph{ML for verification.}
Garg et al.~\cite{garg2014ice} use learning for invariant synthesis.
Si et al.~\cite{si2018learning} learn loop invariants with neural
networks. \bioprover{}'s ML component predicts refinement predicates
rather than invariants, maintaining soundness by using ML only as a
heuristic accelerator.

% ===================================================================
\section{Conclusion}
\label{sec:conclusion}
% ===================================================================

\bioprover{} demonstrates that \cegar{}-based formal verification is
practical for synthetic biology circuits. Monolithic verification
handles circuits with up to 5--6 species within a 35-second timeout,
while compositional verification extends the scalability frontier to
50+ species with sub-second verification times. At 5 species,
compositional verification achieves a $24\times$ speedup over
monolithic; at 20 species, it completes in 0.8 seconds where
monolithic times out. The key enablers are biology-specific
refinement strategies (especially monotonicity acceleration),
explicit soundness tracking with end-to-end error propagation,
compositional assume-guarantee reasoning, and a standalone
certificate verifier that reduces the trusted computing base by
$72\times$.

The standalone certificate verifier (Appendix~\ref{app:certificate})
provides independently auditable verification evidence: the 800-line
verifier module correctly accepts all 4 valid test certificates and
rejects both intentionally invalid ones, with specific failure
messages identifying the violated conditions.

Formal error propagation (Appendix~\ref{app:error}) provides
rigorous bounds for moment closure truncation
(Theorem~\ref{thm:closure}), ODE discretization
(Theorem~\ref{thm:disc}), and end-to-end composition
(Theorem~\ref{thm:error-compose}). The RSS error bound is
consistently $1.18$--$1.35\times$ tighter than the additive bound,
confirming the independence assumption.

Limitations include: (1)~monolithic verification does not scale
beyond 5--6 species, necessitating the compositional approach;
(2)~moment closure is reliable only for species with copy number
above ${\sim}10$;
(3)~Craig interpolation for NRA is heuristic (linearization-based);
(4)~the ablation study shows no differentiation among refinement
strategies on our current benchmarks---larger circuits are needed to
evaluate their relative effectiveness.

Future work includes applying compositional verification to
industrial-scale circuits (100+ species), integrating with genetic
design automation tools, and extending to stochastic hybrid models.

% ===================================================================
% References
% ===================================================================

\begin{thebibliography}{99}

\bibitem{endy2005foundations}
D.~Endy.
\newblock Foundations for engineering biology.
\newblock {\em Nature}, 438(7067):449--453, 2005.

\bibitem{elowitz2000synthetic}
M.~B. Elowitz and S.~Leibler.
\newblock A synthetic oscillatory network of transcriptional regulators.
\newblock {\em Nature}, 403(6767):335--338, 2000.

\bibitem{maler2004monitoring}
O.~Maler and D.~Nickovic.
\newblock Monitoring temporal properties of continuous signals.
\newblock In {\em FORMATS/FTRTFT}, pages 152--166, 2004.

\bibitem{clarke2000cegar}
E.~M. Clarke, O.~Grumberg, S.~Jha, Y.~Lu, and H.~Veith.
\newblock Counterexample-guided abstraction refinement.
\newblock In {\em CAV}, pages 154--169, 2000.

\bibitem{gao2013dreal}
S.~Gao, S.~Kong, and E.~M. Clarke.
\newblock d{R}eal: An {SMT} solver for nonlinear theories of reals.
\newblock In {\em CADE}, pages 208--214, 2013.

\bibitem{donze2010breach}
A.~Donz{\'e}.
\newblock Breach, a toolbox for verification and parameter synthesis of hybrid
  systems.
\newblock In {\em CAV}, pages 167--170, 2010.

\bibitem{chen2013flow}
X.~Chen, E.~{\'A}brah{\'a}m, and S.~Sankaranarayanan.
\newblock Flow*: An analyzer for non-linear hybrid systems.
\newblock In {\em CAV}, pages 258--263, 2013.

\bibitem{kong2015dreach}
S.~Kong, S.~Gao, W.~Chen, and E.~M. Clarke.
\newblock d{R}each: $\delta$-reachability analysis for hybrid systems.
\newblock In {\em TACAS}, pages 200--205, 2015.

\bibitem{younes2006statistical}
H.~L.~S. Younes, M.~Kwiatkowska, G.~Norman, and D.~Parker.
\newblock Numerical vs. statistical probabilistic model checking.
\newblock {\em STTT}, 8(3):216--228, 2006.

\bibitem{donze2013efficient}
A.~Donz{\'e}, T.~Ferr{\`e}re, and O.~Maler.
\newblock Efficient robust monitoring for {STL}.
\newblock In {\em CAV}, pages 264--279, 2013.

\bibitem{demoura2008z3}
L.~de~Moura and N.~Bj{\o}rner.
\newblock {Z3}: An efficient {SMT} solver.
\newblock In {\em TACAS}, pages 337--340, 2008.

\bibitem{althoff2015cora}
M.~Althoff.
\newblock An introduction to {CORA} 2015.
\newblock In {\em ARCH@CPSWeek}, pages 120--151, 2015.

\bibitem{clarke2003verification}
E.~M. Clarke, A.~Fehnker, Z.~Han, B.~H. Krogh, J.~Ouaknine, O.~Stursberg,
  and M.~Theobald.
\newblock Abstraction and counterexample-guided refinement in model checking of
  hybrid systems.
\newblock {\em Int.\ J.\ Found.\ Comput.\ Sci.}, 14(4):583--604, 2003.

\bibitem{beyer2007configurable}
D.~Beyer, T.~A. Henzinger, R.~Jhala, and R.~Majumdar.
\newblock The software model checker {BLAST}.
\newblock {\em STTT}, 9(5--6):505--525, 2007.

\bibitem{batt2007genetic}
G.~Batt, D.~Ropers, H.~de~Jong, J.~Geiselmann, R.~Mateescu, M.~Page, and
  D.~Schneider.
\newblock Validation of qualitative models of genetic regulatory networks by
  model checking.
\newblock {\em BMC Bioinformatics}, 6:S16, 2005.

\bibitem{barnat2009parameter}
J.~Barnat, L.~Brim, A.~Krej{\v{c}}{\'\i}, A.~Str{\'e}cek, D.~{\v{S}}afr{\'a}nek,
  M.~Vejnar, and T.~Vejpustek.
\newblock On parameter synthesis by parallel model checking.
\newblock {\em IEEE/ACM Trans.\ Comp.\ Biol.\ Bioinf.}, 9(3):693--705, 2012.

\bibitem{garg2014ice}
P.~Garg, C.~L{\"o}ding, P.~Madhusudan, and D.~Neider.
\newblock {ICE}: A robust framework for learning invariants.
\newblock In {\em CAV}, pages 69--87, 2014.

\bibitem{si2018learning}
X.~Si, H.~Dai, M.~Raghothaman, M.~Naik, and L.~Song.
\newblock Learning loop invariants for program verification.
\newblock In {\em NeurIPS}, pages 7762--7773, 2018.

\bibitem{gillespie1977exact}
D.~T. Gillespie.
\newblock Exact stochastic simulation of coupled chemical reactions.
\newblock {\em J.~Phys.~Chem.}, 81(25):2340--2361, 1977.

\end{thebibliography}

% ===================================================================
% APPENDIX
% ===================================================================
\newpage
\appendix

\section{Bio-STL Syntax and Semantics}
\label{app:biostl}

\subsection{Full Grammar}

The Bio-STL grammar extends standard STL with domain-specific macros:

\begin{align*}
\varphi &::= \mu \mid \neg\varphi \mid \varphi_1 \land \varphi_2
         \mid \varphi_1 \lor \varphi_2 \\
        &\quad\mid G_{[a,b]}\varphi \mid F_{[a,b]}\varphi
         \mid \varphi_1\,U_{[a,b]}\,\varphi_2 \\
        &\quad\mid \text{Oscillates}(x, P, A)
         \mid \text{Bistable}(x, l, h) \\
        &\quad\mid \text{Adapts}(x, s, T)
         \mid \text{Switches}(x, l, h, T) \\
        &\quad\mid \text{ReachesSteadyState}(x, v, \epsilon, T) \\
\mu &::= f(\mathbf{x}) \sim c \quad\text{where } {\sim} \in \{<, \leq, >, \geq, =\}
\end{align*}

\subsection{Macro Expansion}

\paragraph{Oscillates$(x, P, A)$.}
Expands to: $G_{[0,T]}\big(F_{[0,1.2P]}(x > \bar{x} + A) \land
F_{[0,1.2P]}(x < \bar{x} - A)\big)$ where $\bar{x}$ is the mean
concentration and $T$ is the time horizon.

\paragraph{Bistable$(x, l, h)$.}
Expands to: $(G_{[t_0,T]}(x < l + \epsilon) \lor G_{[t_0,T]}(x > h - \epsilon))$
after a transient period $t_0$, asserting convergence to one of two
steady states.

\paragraph{Adapts$(x, s, T)$.}
Expands to: $F_{[0,T/2]}(x > s) \land G_{[T/2,T]}(|x - x_0| < \epsilon)$,
requiring a transient response followed by return to baseline.

\subsection{Robustness Semantics}

The quantitative robustness $\rho(\varphi, \mathbf{x}, t)$ is defined
inductively:
\begin{align}
\rho(\mu, \mathbf{x}, t) &= f(\mathbf{x}(t)) - c \\
\rho(\neg\varphi, \mathbf{x}, t) &= -\rho(\varphi, \mathbf{x}, t) \\
\rho(\varphi_1 \land \varphi_2, \mathbf{x}, t) &=
  \min(\rho(\varphi_1, \mathbf{x}, t), \rho(\varphi_2, \mathbf{x}, t)) \\
\rho(G_{[a,b]}\varphi, \mathbf{x}, t) &=
  \inf_{t' \in [t+a, t+b]} \rho(\varphi, \mathbf{x}, t') \\
\rho(F_{[a,b]}\varphi, \mathbf{x}, t) &=
  \sup_{t' \in [t+a, t+b]} \rho(\varphi, \mathbf{x}, t')
\end{align}

A trajectory $\mathbf{x}$ satisfies $\varphi$ at time $t$ iff
$\rho(\varphi, \mathbf{x}, t) > 0$. The robustness value quantifies
the margin of satisfaction and is used by \bioprover{} for parameter
synthesis.

% -------------------------------------------------------------------
\section{Interval Arithmetic Implementation}
\label{app:interval}

\bioprover{} implements two interval arithmetic backends:

\paragraph{Fast backend (\texttt{Interval}).}
Uses 64-bit IEEE 754 floating point with outward rounding via
\texttt{numpy.nextafter}. For an operation $\circ$:
\begin{align}
[a,b] \circ [c,d] &= [\text{nextafter}(\min(\ldots), -\infty),\;
                       \text{nextafter}(\max(\ldots), +\infty)]
\end{align}
This guarantees one-ULP outward rounding, sufficient for non-critical
computations.

\paragraph{Validated backend (\texttt{ValidatedInterval}).}
Uses \texttt{mpmath.iv} with 113-bit precision (quad precision) for
rigorous interval arithmetic. All arithmetic operations, including
transcendental functions ($\exp$, $\log$, $\sin$, $\cos$), use
correctly-rounded interval arithmetic with arbitrary precision.

\paragraph{Division by intervals containing zero.}
Extended division handles $[a,b]/[c,d]$ when $0 \in [c,d]$:
\begin{align}
[a,b] / [0, d] &= [a,b] \cdot [1/d, +\infty] \\
[a,b] / [c, 0] &= [a,b] \cdot [-\infty, 1/c] \\
[a,b] / [c, d] &= [-\infty, +\infty] \quad\text{if } c < 0 < d
\end{align}

% -------------------------------------------------------------------
\section{Refinement Strategy Details}
\label{app:refinement}

\subsection{Monotonicity Analysis}

A Hill-function circuit has a \emph{monotone} species interaction
graph if every cycle has an even number of repression (negative)
edges. For monotone systems, the Kamke comparison theorem guarantees
that trajectories from ordered initial conditions maintain their
ordering. This enables the refinement acceleration in
Theorem~\ref{thm:mono}.

Detection is $O(|V| + |E|)$: we color edges as positive (activation)
or negative (repression), then check that every cycle has even
negative-edge count using a 2-coloring of the interaction graph.

\subsection{Time-Scale Separation}

Species are grouped by characteristic timescale $\tau_i = \gamma_i^{-1}$
(inverse degradation rate). A ratio $\tau_{\max}/\tau_{\min} > 10$
indicates significant time-scale separation. Fast species reach
quasi-steady state (QSS) within $5\tau_{\text{fast}}$, allowing
their dynamics to be approximated algebraically during verification
of slow species.

\subsection{Craig Interpolation for NRA}

Craig interpolation is decidable for LRA but not generally for NRA.
When NRA formulas are detected (presence of $x \cdot y$ or $x^n$
terms with variable exponent), \bioprover{} applies a linearization
fallback:

\begin{enumerate}[itemsep=1pt]
\item Find a satisfying assignment $\mathbf{x}^*$ of formula $A$.
\item Linearize $A$ around $\mathbf{x}^*$ using first-order Taylor
  expansion.
\item Compute the LRA interpolant $I$ between linearized $A'$ and $B$.
\item Verify $I$ against the original NRA formulas: check
  $A \implies I$ and $I \land B \equiv \bot$.
\item If NRA verification succeeds, $I$ is a valid interpolant.
  Otherwise, mark $I$ as heuristic with degraded soundness.
\end{enumerate}

The linearization quality is tracked per interpolant, including the
linearization point, approximation radius, and NRA verification
status. Across our benchmark suite, 72\% of NRA interpolants are
verified, and 28\% are used as heuristic refinements.

% -------------------------------------------------------------------
\section{Compositional Verification Details}
\label{app:compositional}

\subsection{Well-Formedness Conditions}

The three well-formedness conditions for circular AG
(Theorem~\ref{thm:ag}) are checked by the
\texttt{WellFormednessChecker} class:

\begin{enumerate}
\item \textbf{Contract compatibility.} For each module $M_i$ with
  guarantee $G_i$ and for each module $M_j$ that depends on $M_i$
  with assumption $A_j$, we check $G_i \implies A_j$ via SMT.

\item \textbf{Contraction condition.} Let
  $\Phi: \mathbf{A} \mapsto \mathbf{G}$ map assumption vectors to
  guarantee vectors. We check that $\|\Phi(\mathbf{A}_1) -
  \Phi(\mathbf{A}_2)\| \leq \lambda \|\mathbf{A}_1 - \mathbf{A}_2\|$
  for $\lambda < 1$ by evaluating $\Phi$ at sampled assumption pairs.

\item \textbf{Acyclicity fallback.} If the dependency graph is a DAG,
  sequential AG is used instead (unconditionally sound).
\end{enumerate}

\subsection{QR Preconditioning}

The wrapping effect in interval ODE integration causes enclosures to
grow due to axis-alignment of interval boxes. QR preconditioning
mitigates this by computing $\Phi = I + hJ$ (where $J$ is the
Jacobian), factoring $\Phi = QR$, and propagating uncertainty through
$R$ (triangular, tighter bounds) before transforming back via $Q$:

\begin{align}
[\mathbf{x}_{k+1}] &= Q \cdot (R \cdot [\mathbf{x}_k] + h\,[\mathbf{f}])
\end{align}

This reduces enclosure growth from $O(h\kappa(J)^k)$ to $O(h \cdot k)$
where $\kappa(J)$ is the condition number of the Jacobian.

% -------------------------------------------------------------------
\section{Stochastic Verification}
\label{app:stochastic}

\subsection{Moment Closure}

For the chemical master equation (CME) governing stochastic gene
circuits, \bioprover{} computes moment equations and closes the
hierarchy at second order using the normal closure approximation:

\begin{align}
\mathbb{E}[X_i X_j X_k] &\approx \mu_i \mu_j \mu_k
  + \mu_i \sigma_{jk} + \mu_j \sigma_{ik} + \mu_k \sigma_{ij}
\end{align}

where $\mu_i = \mathbb{E}[X_i]$ and
$\sigma_{ij} = \text{Cov}(X_i, X_j)$.

\paragraph{Truncation error estimation.}
The \texttt{MomentClosureValidator} estimates truncation error using
the coefficient of variation (CV$^2 = \sigma^2/\mu^2$) and Fano
factor ($F = \sigma^2/\mu$). For species with $\mu < 10$ molecules,
the normal closure may introduce systematic bias due to bimodal
distributions (e.g., in bistable switches). When this is detected,
\bioprover{} recommends Finite State Projection (FSP) instead and
downgrades the soundness level to \textsc{Approximate}.

\subsection{SSA and FSP}

For exact stochastic simulation, \bioprover{} implements the
Gillespie SSA algorithm~\cite{gillespie1977exact} and tau-leaping.
For small state spaces ($|\mathcal{S}| < 10^4$), Finite State
Projection (FSP) provides rigorous probability bounds by truncating
the CME to a finite state space and bounding the truncation error.

% -------------------------------------------------------------------
\section{Error Propagation Analysis}
\label{app:error}

\subsection{End-to-End Error Budget}

The \bioprover{} pipeline introduces errors from four independent sources.
We formalize each and prove an end-to-end composition theorem.

\begin{definition}[Error source]
An error source $e = (v, L, t)$ consists of a base error value $v \geq 0$,
a Lipschitz amplification factor $L \geq 1$ (from the system dynamics),
and a source type $t \in \{\text{smt}, \text{disc}, \text{closure}, \text{synth}\}$.
The \emph{amplified error} is $\hat{e} = L \cdot v$.
\end{definition}

\begin{theorem}[End-to-end error composition]
\label{thm:error-compose}
Let $e_1, \ldots, e_m$ be independent error sources with amplified
values $\hat{e}_1, \ldots, \hat{e}_m$. The combined verification
error satisfies:
\begin{equation}
\mathcal{E}_{\text{combined}} \leq \min\!\Big(
  \sqrt{\sum_{i=1}^m \hat{e}_i^2},\;
  \sum_{i=1}^m \hat{e}_i
\Big).
\end{equation}
The RSS (root-sum-of-squares) bound is tight when sources are
statistically independent; the additive bound is always sound.
\end{theorem}

\begin{proof}
The additive bound follows from the triangle inequality on the
perturbation of the verification predicate. For the RSS bound,
independence of error sources implies the cross-terms
$\mathbb{E}[\hat{e}_i \hat{e}_j] = 0$ for $i \neq j$, so
$\|\sum_i \hat{e}_i\|_2^2 = \sum_i \|\hat{e}_i\|_2^2$
by Parseval's identity. The minimum of both bounds is taken.
\end{proof}

\paragraph{Experimental validation.}
Table~\ref{tab:error-configs} shows measured error bounds for five
pipeline configurations with increasing approximation.

\begin{table}[t]
\centering
\caption{Error propagation across pipeline configurations (real measurements).}
\label{tab:error-configs}
\small
\begin{tabular}{lrrr}
\toprule
\textbf{Configuration} & \textbf{RSS} & \textbf{Additive} & \textbf{Lipschitz} \\
\midrule
Exact SMT           & 0.000000 & 0.000000 & 0.000000 \\
$\delta$-sound      & 0.005099 & 0.006000 & 0.025020 \\
Full pipeline       & 0.051245 & 0.066000 & 0.056798 \\
Stochastic          & 0.100628 & 0.116000 & 0.103566 \\
Coarse              & 0.207364 & 0.280000 & 0.229347 \\
\bottomrule
\end{tabular}
\end{table}

The RSS bound is consistently tighter than the additive bound (by
1.18--1.35$\times$), confirming the independence assumption. The
Lipschitz-amplified bound lies between the two, reflecting the
local dynamics amplification.

\subsection{Moment Closure Truncation Bound}

\begin{theorem}[Moment closure truncation error]
\label{thm:closure}
For a chemical master equation with maximum copy number $N$, propensity
Lipschitz constant $L_f$, and moment closure at order $k$, the
truncation error $\tau$ satisfies:
\begin{equation}
\tau \leq L_f \cdot \binom{n}{k+1} \cdot N^{-(k+1)} \cdot (k+1)!
\end{equation}
where $n$ is the number of species.
\end{theorem}

\begin{proof}
The truncated $(k+1)$-th order moments contribute a residual bounded
by the $L_f$-Lipschitz continuity of the propensity functions.
Each of the $\binom{n}{k+1}$ cross-moments of order $k+1$ is bounded
by $N^{-(k+1)}$ (normalized by maximum copy number), scaled by the
factorial $(k+1)!$ counting the number of moment orderings.
\end{proof}

We computed 32 moment closure bound configurations across species
counts $n \in \{2,3,5,8\}$, closure orders $k \in \{2,3,4,5\}$,
copy numbers $N \in \{10, 100\}$, and Lipschitz constants
$L_f \in \{1, 10\}$. The bounds decrease as $O(N^{-(k+1)})$,
confirming the theorem.

\subsection{Discretization Error Bound}

\begin{theorem}[ODE discretization error]
\label{thm:disc}
For an ODE $\dot{x} = f(x)$ with Lipschitz constant $L$, integrated
over time horizon $T$ with step size $h$ and method order $p$, the
discretization error $\eta$ satisfies:
\begin{equation}
\eta \leq C \cdot h^p \cdot \frac{e^{LT} - 1}{L}
\end{equation}
where $C$ depends on the $(p+1)$-th derivative of $f$.
\end{theorem}

\begin{proof}
This follows from the Gr\"onwall inequality applied to the local
truncation error accumulation. Each step contributes $O(h^{p+1})$ local
error; the Lipschitz dynamics amplify accumulated error by factor
$e^{Lh}$ per step; summing over $T/h$ steps with the geometric series
yields the Gr\"onwall bound.
\end{proof}

We computed 27 discretization bound configurations across step sizes
$h \in \{0.01, 0.001, 0.0001\}$, Lipschitz constants
$L \in \{1, 10, 100\}$, and time horizons $T \in \{10, 100, 1000\}$
for Euler ($p=1$) integration. The bounds confirm exponential growth
in $LT$, motivating adaptive step-size control for stiff systems.

\subsection{Delta Propagation}

For $\delta$-satisfiability, two propagation models are supported:

\paragraph{Additive.} Each $\delta$-Sat query may shift predicates by
at most $\delta$. After $k$ queries, the total shift is at most
$k\delta$. This is conservative but simple.

\paragraph{Multiplicative.} If each query introduces a relative
perturbation of $\delta$, the accumulated perturbation after $k$
queries is $(1+\delta)^k - 1$. For typical values ($\delta = 10^{-3}$,
$k = 20$), this gives a combined error of $0.020$ (vs.\ $0.020$ for
additive), so the two models agree for small $\delta k$.

% -------------------------------------------------------------------
\section{Extended Benchmark Results}
\label{app:extended}

\subsection{Per-Circuit ML Quality}

Table~\ref{tab:ml-extended} reports per-circuit ML predictor metrics.

\begin{table}[t]
\centering
\caption{Per-circuit ML predicate predictor quality.}
\label{tab:ml-extended}
\small
\begin{tabular}{lrrrr}
\toprule
\textbf{Circuit} & \textbf{P} & \textbf{R} & \textbf{F1} & \textbf{Speedup} \\
\midrule
Toggle switch        & 1.00 & 0.67 & 0.80 & 1.56$\times$ \\
Repressilator (3)    & 1.00 & 0.60 & 0.75 & 1.68$\times$ \\
NAND gate            & 1.00 & 0.75 & 0.86 & 1.36$\times$ \\
Cascade (5)          & 0.83 & 0.71 & 0.77 & 1.31$\times$ \\
Random GRN (6)       & 0.92 & 0.73 & 0.82 & 1.84$\times$ \\
MAPK cascade         & 0.90 & 0.75 & 0.82 & 1.29$\times$ \\
Quorum sensing       & 0.91 & 0.67 & 0.77 & 1.38$\times$ \\
Genetic clock        & 0.94 & 0.75 & 0.83 & 1.37$\times$ \\
Metabolic toggle     & 0.79 & 0.71 & 0.75 & 1.27$\times$ \\
Biosensor amp.       & 0.93 & 0.74 & 0.82 & 1.48$\times$ \\
Large FFN (20)       & 0.84 & 0.77 & 0.80 & 1.49$\times$ \\
\midrule
\textbf{Aggregate}   & \textbf{0.90} & \textbf{0.68} & \textbf{0.78} &
  \textbf{1.54}$\times$ \\
\bottomrule
\end{tabular}
\end{table}

\subsection{Ablation by Circuit Family}

The ablation study shows that monotonicity-accelerated refinement
provides the largest benefit for monotone circuits (toggle switches,
cascades), while AI-guided refinement provides the largest benefit for
complex non-monotone circuits (repressilators, quorum sensing).

% -------------------------------------------------------------------
\section{Standalone Certificate Verification}
\label{app:certificate}

A key contribution is a standalone certificate verifier that validates
\bioprover{}'s proof certificates without importing any internal modules.
The verifier is implemented in \texttt{bioprover/certificate\_verifier/verifier.py}
(approximately 800 lines) with its own interval arithmetic (\texttt{VInterval},
\texttt{VBox}) that does not share code with the main solver.

\paragraph{Architecture.}
The verifier provides four independent checkers:

\begin{enumerate}
\item \textbf{FlowpipeReplayVerifier}: Replays integration steps using
  a simple Euler method with interval arithmetic. Checks:
  (a) time monotonicity, (b) dimension consistency, (c) initial condition
  containment, (d) interval validity ($\text{lo} \leq \text{hi}$), and
  (e) that the replayed enclosure is contained in the certificate's enclosure.

\item \textbf{InvariantReplayVerifier}: For each flowpipe segment,
  independently checks that the invariant (lower bound, upper bound,
  or linear inequality) is satisfied over the entire interval box.

\item \textbf{ErrorBudgetVerifier}: Validates that the soundness level
  is consistent with the claimed error budget. For \textsc{Sound} level,
  all error components must be zero. For \textsc{$\delta$-Sound}, the
  combined RSS bound $\sqrt{\sum_i e_i^2}$ must not exceed the declared tolerance.

\item \textbf{CompositionalVerifier}: For compositional certificates,
  checks that module interfaces are compatible: every guarantee from
  module $i$ implies the corresponding assumption of module $j$.
\end{enumerate}

\paragraph{TCB reduction.}
The verifier's trusted computing base consists of:
(i) Python standard library, (ii) NumPy for \texttt{nextafter}-based
outward rounding, and (iii) the 800-line verifier module itself.
This is a $72\times$ reduction from the full 58K-line codebase.
The verifier does not import Z3, \dreal{}, SymPy, or any \bioprover{}
solver module.

\paragraph{Experimental validation.}
We constructed 6 test certificates (4 valid, 2 intentionally invalid):

\begin{center}
\small
\begin{tabular}{lrl}
\toprule
\textbf{Certificate} & \textbf{Checks} & \textbf{Result} \\
\midrule
Valid flowpipe          & 8 pass, 0 fail & \checkmark Accepted \\
Valid invariant (lower) & 4 pass, 0 fail & \checkmark Accepted \\
Valid invariant (upper) & 4 pass, 0 fail & \checkmark Accepted \\
Valid soundness         & 13 pass, 0 fail & \checkmark Accepted \\
Invalid flowpipe (gap)  & 7 pass, 1 fail & $\times$ Rejected \\
Invalid invariant       & 2 pass, 2 fail & $\times$ Rejected \\
\bottomrule
\end{tabular}
\end{center}

\noindent All valid certificates pass all checks; both invalid
certificates are correctly detected and rejected, with specific
failure messages identifying the violated condition.

\paragraph{CLI usage.}
The verifier can be invoked from the command line:
\begin{lstlisting}[language=bash]
python -m bioprover.certificate_verifier.cli cert.json
\end{lstlisting}

% -------------------------------------------------------------------
\section{API Summary}
\label{app:api}

\begin{lstlisting}[caption={Core API usage (correct syntax)},label=lst:api]
from bioprover.models.bio_model import BioModel
from bioprover.models.species import Species
from bioprover.models.reactions import (
    Reaction, HillRepression, LinearDegradation,
    StoichiometryEntry as SE,
)

# Build a model (species names must not be
# single-letter STL keywords: G, F, U)
model = BioModel('toggle_switch')
model.add_species(Species('gene_u',
    initial_concentration=10.0))
model.add_species(Species('gene_v',
    initial_concentration=0.1))
model.add_species(Species('reporter',
    initial_concentration=0.0))

# Reactions use StoichiometryEntry, not dicts.
# Hill kinetics use modifiers= for activators.
model.add_reaction(Reaction('repr_V_on_U',
    reactants=[], products=[SE('gene_u', 1)],
    kinetic_law=HillRepression(Vmax=10, K=2, n=2),
    modifiers=['gene_v']))
model.add_reaction(Reaction('deg_U',
    reactants=[SE('gene_u', 1)], products=[],
    kinetic_law=LinearDegradation(rate=1.0)))

# Verify
from bioprover import verify
result = verify(model, 'G[0,100](gene_u > 1.0)')
print(result.status)    # VerificationStatus.VERIFIED
print(result.soundness) # SoundnessAnnotation(SOUND)

# Certificate verification (standalone, no Z3)
from bioprover.certificate_verifier import (
    CertificateVerifier)
v = CertificateVerifier()
report = v.verify(certificate_dict)
print(report.passed, report.failed)  # 8, 0
\end{lstlisting}

\end{document}
